\section{Conclusiones}

La aplicación de K-Means sobre las características de audio del dataset
de Spotify permitió identificar [K] clusters musicales con perfiles
diferenciados. Las principales conclusiones son:

\begin{enumerate}
  \item La normalización es crítica: sin StandardScaler, las variables
        con mayor rango (tempo, loudness) dominarían la distancia euclidiana.
  \item PCA con 2 componentes explica [X]\% de la varianza, suficiente
        para visualización pero no para reconstrucción fiel.
  \item Los clusters identificados corresponden parcialmente a géneros
        musicales reales, validando la coherencia del agrupamiento.
  \item [Completar con hallazgo específico del análisis]
\end{enumerate}

Como trabajo futuro se propone explorar clustering jerárquico para
análisis de subgéneros, y evaluar modelos de mezcla gaussiana (GMM)
como alternativa probabilística a K-Means.

\section*{Agradecimientos}
Los autores agradecen al docente Aimer Rivera Centeno de la Universidad
Popular del Cesar por la orientación académica en este proyecto.
